\documentclass[12pt]{article}
\usepackage[utf8]{inputenc}
\usepackage{amsmath, amssymb, amsthm}
\usepackage{geometry}
\usepackage{xcolor}
\usepackage[most]{tcolorbox}
\usepackage{enumitem}
\usepackage{fancyhdr}

% Page setup
\geometry{margin=1in}
\setlength{\headheight}{14.5pt}
\pagestyle{fancy}
\fancyhf{}
\rhead{AMC12 Problem Analysis}
\lhead{\today}
\cfoot{\thepage}

% Color definitions
\definecolor{problemconfig}{RGB}{230, 242, 255}    % Light blue
\definecolor{strategic}{RGB}{255, 230, 242}       % Light pink
\definecolor{tactical}{RGB}{242, 255, 230}        % Light green
\definecolor{techniques}{RGB}{255, 242, 230}      % Light yellow
\definecolor{verification}{RGB}{255, 230, 230}    % Light red
\definecolor{solution}{RGB}{230, 255, 255}        % Light cyan
\definecolor{competition}{RGB}{242, 242, 255}     % Light purple
\definecolor{gold}{RGB}{218, 165, 32}

% Box styles
\newtcolorbox{problemconfigbox}{
    colback=problemconfig,
    colframe=blue,
    title=Problem Configuration Analysis,
    fonttitle=\bfseries,
    arc=3pt,
    boxrule=1pt,
    breakable
}

\newtcolorbox{strategicbox}{
    colback=strategic,
    colframe=purple,
    title=Strategic Framework Application,
    fonttitle=\bfseries,
    arc=3pt,
    boxrule=1pt,
    breakable
}

\newtcolorbox{tacticalbox}{
    colback=tactical,
    colframe=green,
    title=Tactical Methodology Selection,
    fonttitle=\bfseries,
    arc=3pt,
    boxrule=1pt,
    breakable
}

\newtcolorbox{techniquesbox}{
    colback=techniques,
    colframe=orange,
    title=Conversion Techniques Application,
    fonttitle=\bfseries,
    arc=3pt,
    boxrule=1pt,
    breakable
}

\newtcolorbox{verificationbox}{
    colback=verification,
    colframe=red,
    title=Verification Strategy Design,
    fonttitle=\bfseries,
    arc=3pt,
    boxrule=1pt,
    breakable
}

\newtcolorbox{solutionbox}{
    colback=solution,
    colframe=cyan,
    title=Solution Roadmap,
    fonttitle=\bfseries,
    arc=3pt,
    boxrule=1pt,
    breakable
}

\newtcolorbox{competitionbox}{
    colback=competition,
    colframe=violet,
    title=Competition Strategy Integration,
    fonttitle=\bfseries,
    arc=3pt,
    boxrule=1pt,
    breakable
}

% Highlight boxes
\newtcolorbox{keyinsight}{
    colback=gold!10,
    colframe=gold,
    title=Key Insight,
    fonttitle=\bfseries,
    arc=3pt,
    boxrule=1.5pt,
    breakable
}

\newtcolorbox{warningbox}{
    colback=red!10,
    colframe=red,
    title=Warning,
    fonttitle=\bfseries,
    arc=3pt,
    boxrule=1.5pt,
    breakable
}

\newtcolorbox{protip}{
    colback=green!10,
    colframe=green!70!black,
    title=Pro Tip,
    fonttitle=\bfseries,
    arc=3pt,
    boxrule=1.5pt,
    breakable
}

\title{\textbf{AMC 12A 2024 Problem 24 Analysis}\\
\large{Comprehensive Strategic Framework Application}}
\author{Problem Analysis Framework v2.1}
\date{\today}

\begin{document}

\maketitle

\section*{Problem Statement}

A \textit{disphenoid} is a tetrahedron whose triangular faces are congruent to one another. What is the least total surface area of a disphenoid whose faces are scalene triangles with integer side lengths?

\begin{align*}
\textbf{(A) } \sqrt{3} \qquad \textbf{(B) } 3\sqrt{15} \qquad \textbf{(C) } 15 \qquad \textbf{(D) } 15\sqrt{7} \qquad \textbf{(E) } 24\sqrt{6}
\end{align*}

\section{Problem Configuration Analysis}

\begin{problemconfigbox}

\subsection*{A. Problem Classification}

\textbf{Problem Type:} To Find (computational optimization problem)

\textbf{Difficulty Band:} Late-band (Problem 24 of 25)
\begin{itemize}
    \item Expected time: 5-8 minutes
    \item Difficulty level: Very High (96\% position)
    \item Requires deep geometric understanding and optimization
\end{itemize}

\textbf{Competition Context:} AMC 12A (75 minutes, 25 problems, multiple choice)

\textbf{Mathematical Domain:} 
\begin{itemize}
    \item Primary: 3D Geometry (tetrahedra, disphenoids)
    \item Secondary: Number Theory (integer side lengths)
    \item Tertiary: Optimization (minimization problem)
    \item Quaternary: Triangle Geometry (Heron's Formula, acute triangles)
\end{itemize}

\subsection*{B. Problem Structure Identification}

\textbf{Given Information:}
\begin{itemize}
    \item Definition: Disphenoid = tetrahedron with 4 congruent triangular faces
    \item Constraint: Faces must be scalene triangles
    \item Constraint: Side lengths must be integers
    \item Objective: Minimize total surface area
\end{itemize}

\textbf{Constraints:}
\begin{itemize}
    \item All four faces are congruent (same shape and size)
    \item Faces are scalene: $a \neq b \neq c \neq a$
    \item Side lengths are positive integers: $a, b, c \in \mathbb{Z}^+$
    \item Must form valid triangle: Triangle inequality must hold
    \item Must form valid disphenoid: Geometric realizability condition
\end{itemize}

\textbf{Objective:} Find the minimum value of $4A$ where $A$ is the area of one face.

\textbf{Hidden Assumptions:}
\begin{itemize}
    \item The disphenoid must actually be constructible (geometric realizability)
    \item For a disphenoid with scalene faces, opposite edges must be equal
    \item The faces must form an acute triangle (discovered through analysis)
    \item Smaller integers don't necessarily mean smaller area
\end{itemize}

\textbf{Answer Choice Leverage:}
\begin{itemize}
    \item All choices involve radicals, suggesting Heron's formula
    \item Choice (C) $15$ is suspiciously simple (likely trap for non-scalene)
    \item Choice (D) $15\sqrt{7}$ has form matching 4-5-6 triangle calculation
    \item Answer structure suggests looking for triangles with semi-integer semiperimeter
\end{itemize}

\subsection*{C. Strategic Orientation}

\textbf{Hypothesis:} A disphenoid exists with scalene integer-sided faces.

\textbf{Conclusion:} Determine the minimum possible surface area.

\textbf{Notation:}
\begin{itemize}
    \item Let $a, b, c$ be the side lengths of one triangular face
    \item Let $A$ be the area of one face
    \item Total surface area = $4A$
    \item Let $s = \frac{a+b+c}{2}$ be the semiperimeter
\end{itemize}

\textbf{Intuitive Approach:} 
\begin{itemize}
    \item Start with smallest possible scalene integer triangles
    \item Check if they can form a valid disphenoid
    \item Calculate area using Heron's formula
    \item Work systematically through candidates
\end{itemize}

\textbf{Cross-Domain Potential:}
\begin{itemize}
    \item Algebraic: Disphenoid can be inscribed in a rectangular box
    \item Analytic: System of equations for box dimensions
    \item Number-theoretic: Integer constraints create discrete optimization
\end{itemize}

\end{problemconfigbox}

\begin{keyinsight}
\textbf{Critical Discovery:} A disphenoid with scalene faces must have its faces as \textit{acute} triangles. This is because when a disphenoid is inscribed in a rectangular box with dimensions $p, q, r$, we have:
\begin{align*}
p^2 &= \frac{a^2 + b^2 - c^2}{2} \\
q^2 &= \frac{a^2 - b^2 + c^2}{2} \\
r^2 &= \frac{-a^2 + b^2 + c^2}{2}
\end{align*}
For $p, q, r$ to be real and positive (required for a physical box), all three expressions must be positive, which by the law of cosines means all angles of the triangle must be acute.
\end{keyinsight}

\section{Strategic Framework Application}

\begin{strategicbox}

\subsection*{A. Psychological Strategies}

\textbf{Mental Toughness:}
\begin{itemize}
    \item This is Problem 24 - expect multiple failed attempts
    \item Don't panic if first few triangles don't work
    \item The acute triangle constraint is non-obvious and must be discovered
    \item Stay systematic rather than random guessing
\end{itemize}

\textbf{Creativity Cultivation:}
\begin{itemize}
    \item Consider the geometric constraint: what makes a disphenoid special?
    \item Explore the relationship between disphenoid and rectangular box
    \item Think about why opposite edges must be equal
    \item Visualize the 3D structure if possible
\end{itemize}

\textbf{Iterative Refinement:}
\begin{itemize}
    \item Test candidate triangles: 2-3-4, 3-4-5, 4-5-6, etc.
    \item For each failure, understand \textit{why} it failed
    \item Develop the acute triangle criterion through testing
    \item Refine search strategy based on patterns
\end{itemize}

\subsection*{B. Investigation Strategies}

\textbf{Startup Orientation:}
\begin{itemize}
    \item Understand the definition of disphenoid
    \item Recognize this is geometric optimization with integer constraints
    \item Note that ``least'' means we need to check small cases systematically
    \item Identify that answer choices give hints about expected form
\end{itemize}

\textbf{Get Your Hands Dirty:}
\begin{itemize}
    \item List smallest scalene triangles: (2,3,4), (2,4,5), (3,4,5), (4,5,6)
    \item Test each for valid disphenoid formation
    \item Calculate area for valid candidates using Heron's formula
    \item Pattern: smaller sides don't always mean smaller area
\end{itemize}

\textbf{Penultimate Step:}
\begin{itemize}
    \item Final step: Multiply area of one face by 4
    \item Before that: Calculate area using Heron's formula
    \item Before that: Identify the smallest valid scalene acute triangle
    \item Before that: Understand the acute triangle requirement
\end{itemize}

\textbf{Wishful Thinking:}
\begin{itemize}
    \item Wish: If only all scalene triangles could form disphenoids
    \item Reality: Only acute triangles work
    \item Wish: If only (2,3,4) worked (smallest scalene)
    \item Reality: Must use (4,5,6) as smallest acute scalene
\end{itemize}

\textbf{Make It Easier:}
\begin{itemize}
    \item Simplify: First find smallest scalene triangle (2,3,4)
    \item Discover: It's obtuse, doesn't work
    \item Next: Find smallest \textit{acute} scalene triangle
    \item Test systematically: (3,4,5) is right, (4,5,6) is acute!
\end{itemize}

\textbf{Conjecture via Small Cases:}
\begin{itemize}
    \item Test (2,3,4): $4 + 9 = 13 < 16$ (obtuse)
    \item Test (2,4,5): $4 + 16 = 20 < 25$ (obtuse)
    \item Test (3,4,5): $9 + 16 = 25$ (right angle, not scalene in angle)
    \item Test (4,5,6): $16 + 25 = 41 > 36$ (acute!) $\checkmark$
    \item Pattern: Need $a^2 + b^2 > c^2$ for all pairings
\end{itemize}

\textbf{Visualization:}
\begin{itemize}
    \item Imagine a rectangular box with vertices at coordinates
    \item Disphenoid vertices at alternating corners: $(0,0,0)$, $(p,q,0)$, $(0,q,r)$, $(p,0,r)$
    \item Face diagonals have lengths $\sqrt{p^2+q^2}$, $\sqrt{p^2+r^2}$, $\sqrt{q^2+r^2}$
    \item This visualization reveals the box-embedding structure
\end{itemize}

\subsection*{C. Late-Band Strategic Playbook}

\textbf{Answer-Choice Leverage:}
\begin{itemize}
    \item Eliminate (A) $\sqrt{3}$: Too small for integer-sided triangle
    \item Eliminate (C) $15$: Integer answer unlikely for scalene requirement
    \item Focus on (B), (D), (E) with radical forms
    \item Choice (D) $15\sqrt{7}$ matches structure $4 \cdot \frac{15}{4}\sqrt{7}$
\end{itemize}

\textbf{Make It Easier Application:}
\begin{itemize}
    \item Replace ``disphenoid'' with simpler question: ``What scalene triangles work?''
    \item Replace ``minimum area'' with ``smallest integer sides''
    \item Build back complexity: Add acute constraint
    \item Build back complexity: Calculate actual area
\end{itemize}

\textbf{Penultimate Step Focus:}
\begin{itemize}
    \item Working backwards: Need $15\sqrt{7}$ total area
    \item Therefore: Need $\frac{15\sqrt{7}}{4}$ per face
    \item Therefore: Need triangle with this specific area
    \item Check: Does (4,5,6) give this? Yes!
\end{itemize}

\subsection*{D. Argument Construction Strategy}

\textbf{Direct Proof Approach:}
\begin{enumerate}
    \item Establish that disphenoid faces must be acute triangles
    \item Identify smallest scalene acute triangle with integer sides
    \item Calculate its area using Heron's formula
    \item Multiply by 4 for total surface area
\end{enumerate}

\textbf{Key Lemma:} For a disphenoid with opposite edges equal (essential property), if faces have sides $a, b, c$, they must satisfy the acute triangle condition: $a^2 + b^2 > c^2$ for all permutations.

\end{strategicbox}

\begin{warningbox}
\textbf{Common Mistake:} Assuming (2,3,4) or (3,4,5) work without checking the acute triangle requirement. Many students calculate area for these triangles and get wrong answers. The geometric realizability constraint is the key difficulty in this problem.
\end{warningbox}

\section{Tactical Methodology Selection}

\begin{tacticalbox}

\subsection*{A. Core Mathematical Tactics}

\textbf{Extreme Principle:}
\begin{itemize}
    \item Applied: Seeking minimum surface area
    \item Strategy: Test smallest possible integer triangles first
    \item Refinement: Add acute constraint to narrow search
    \item Success: (4,5,6) is extremal under constraints
\end{itemize}

\textbf{Systematic Case Analysis:}
\begin{itemize}
    \item Enumerate small scalene triangles: (2,3,4), (2,4,5), (3,4,5), (4,5,6), (3,5,7)
    \item For each, check: Is it acute?
    \item For acute ones, calculate area
    \item Select minimum area
\end{itemize}

\textbf{Invariant Principle:}
\begin{itemize}
    \item Key invariant: In a disphenoid, opposite edges are equal
    \item This forces the box-embedding structure
    \item This creates the acute triangle requirement
    \item Understanding this invariant is crucial
\end{itemize}

\textbf{Algebraic Manipulation:}
\begin{itemize}
    \item Use Heron's formula: $A = \sqrt{s(s-a)(s-b)(s-c)}$
    \item For (4,5,6): $s = \frac{15}{2}$
    \item Calculate: $A = \sqrt{\frac{15}{2} \cdot \frac{7}{2} \cdot \frac{5}{2} \cdot \frac{3}{2}}$
    \item Simplify: $A = \frac{15\sqrt{7}}{4}$
\end{itemize}

\subsection*{B. Domain-Specific Tactics}

\textbf{Geometry Tactics:}
\begin{itemize}
    \item \textbf{3D Visualization:} Box embedding of disphenoid
    \item \textbf{Constraint Analysis:} Opposite edges equal
    \item \textbf{Pythagorean Relationships:} Box dimensions from triangle sides
    \item \textbf{Triangle Classification:} Acute vs. right vs. obtuse
\end{itemize}

\textbf{Number Theory Tactics:}
\begin{itemize}
    \item \textbf{Integer Constraints:} Discrete optimization space
    \item \textbf{Small Case Enumeration:} Systematic testing
    \item \textbf{Divisibility:} Looking for perfect square factors
\end{itemize}

\textbf{Optimization Tactics:}
\begin{itemize}
    \item \textbf{Greedy Approach:} Start with smallest integers
    \item \textbf{Constraint Satisfaction:} Filter by acute requirement
    \item \textbf{Monotonicity:} Generally smaller sides $\to$ smaller area
\end{itemize}

\subsection*{C. Crossover Techniques}

\textbf{Geometry $\leftrightarrow$ Algebra:}
\begin{itemize}
    \item Translate: 3D disphenoid $\to$ system of equations for box
    \item Solve: Box dimensions in terms of triangle sides
    \item Discover: Acute triangle requirement from positive dimensions
\end{itemize}

\textbf{Optimization $\leftrightarrow$ Number Theory:}
\begin{itemize}
    \item Combine: Minimization with integer constraints
    \item Result: Finite search space over small integers
    \item Strategy: Enumerate and test
\end{itemize}

\end{tacticalbox}

\begin{protip}
\textbf{Time-Saving Recognition:} If you recognize that a disphenoid can be inscribed in a rectangular box with face diagonals as edges, you can immediately write down the system:
\begin{align*}
p^2 + q^2 &= a^2 \\
p^2 + r^2 &= b^2 \\
q^2 + r^2 &= c^2
\end{align*}
Solving this system shows that $p^2, q^2, r^2 \propto$ (cosines of angles), revealing the acute requirement instantly. This saves 2-3 minutes of exploration.
\end{protip}

\section{Conversion Techniques Application}

\begin{techniquesbox}

\subsection*{A. Recognition Index}

\textbf{Primary Technique: Algebraic Translation}
\begin{itemize}
    \item Recognition trigger: ``Disphenoid'' is unfamiliar geometric object
    \item Strategy: Convert to known algebraic framework
    \item Approach: Box-embedding with coordinate geometry
    \item Result: System of Pythagorean equations
\end{itemize}

\textbf{Secondary Technique: Constraint Reformulation}
\begin{itemize}
    \item Recognition trigger: ``Congruent faces'' constraint
    \item Strategy: What does this imply about edge relationships?
    \item Discovery: Opposite edges must be equal
    \item Application: Reduces search space significantly
\end{itemize}

\textbf{Tertiary Technique: Small Case Analysis}
\begin{itemize}
    \item Recognition trigger: ``Integer side lengths'' + ``least''
    \item Strategy: Exhaustive search over small integers
    \item Method: Systematic enumeration with filtering
    \item Efficiency: Very few cases need checking
\end{itemize}

\subsection*{B. Technique Selection}

\textbf{Why Algebraic Translation?}
\begin{itemize}
    \item Disphenoid definition is geometric but abstract
    \item Box-embedding makes it concrete and calculable
    \item Provides necessary and sufficient conditions
    \item Reveals hidden acute triangle constraint
\end{itemize}

\textbf{Why Heron's Formula?}
\begin{itemize}
    \item Given: Three integer side lengths
    \item Need: Area of triangle
    \item Heron's formula is the standard tool
    \item Alternative (base-height) requires more geometry
\end{itemize}

\textbf{Why Systematic Enumeration?}
\begin{itemize}
    \item Search space is small (only need to check up to 4-5-6)
    \item Each check is quick (triangle inequality + acute test)
    \item Guarantees finding global minimum
    \item More reliable than trying to optimize analytically
\end{itemize}

\subsection*{C. Integration Strategy}

\textbf{Step 1: Geometric Understanding}
\begin{itemize}
    \item Understand disphenoid structure
    \item Recognize box-embedding property
    \item Derive acute triangle requirement
\end{itemize}

\textbf{Step 2: Constraint Application}
\begin{itemize}
    \item List scalene integer triangles
    \item Filter by acute condition
    \item Identify (4,5,6) as smallest
\end{itemize}

\textbf{Step 3: Calculation}
\begin{itemize}
    \item Apply Heron's formula to (4,5,6)
    \item Calculate $A = \frac{15\sqrt{7}}{4}$
    \item Multiply by 4: Total = $15\sqrt{7}$
\end{itemize}

\subsection*{D. Conflict Resolution}

\textbf{Potential Conflict 1: Definition vs. Intuition}
\begin{itemize}
    \item Conflict: Intuition says ``smallest sides = smallest area''
    \item Resolution: Not all triangles form valid disphenoids
    \item Lesson: Geometric constraints override simple optimization
\end{itemize}

\textbf{Potential Conflict 2: Scalene vs. Acute}
\begin{itemize}
    \item Conflict: (3,4,5) is scalene but right-angled
    \item Resolution: Right angle fails realizability (one box dimension = 0)
    \item Lesson: Must be strictly acute, not just non-obtuse
\end{itemize}

\textbf{Potential Conflict 3: Multiple Choice Distraction}
\begin{itemize}
    \item Conflict: Choice (C) = 15 looks clean
    \item Resolution: This is $4 \times \frac{15}{4}$ without the $\sqrt{7}$
    \item Likely trap: Area of (3,4,5) without checking validity
\end{itemize}

\end{techniquesbox}

\section{Verification Strategy Design}

\begin{verificationbox}

\subsection*{A. Verification Level Selection}

\textbf{Selected Level: Level 5 - Targeted Deep Check (2-3 minutes)}

\textbf{Decision Tree Analysis:}
\begin{itemize}
    \item Problem difficulty: Late-band (P24) $\to$ High stakes
    \item Personal confidence: Novel geometric object $\to$ Medium confidence
    \item Time remaining: Assume 5-10 minutes left $\to$ Can afford verification
    \item Answer type: Multiple choice with calculation $\to$ Need accuracy
    \item Recommendation: Level 5 (Targeted Deep Check)
\end{itemize}

\subsection*{B. Specialized Verification Methods}

\textbf{Geometry-Specific Checks:}

\textbf{Check 1: Acute Triangle Verification}
\begin{itemize}
    \item Verify $4^2 + 5^2 > 6^2$: $16 + 25 = 41 > 36$ $\checkmark$
    \item Verify $4^2 + 6^2 > 5^2$: $16 + 36 = 52 > 25$ $\checkmark$
    \item Verify $5^2 + 6^2 > 4^2$: $25 + 36 = 61 > 16$ $\checkmark$
    \item Conclusion: (4,5,6) is indeed acute
\end{itemize}

\textbf{Check 2: Triangle Inequality}
\begin{itemize}
    \item $4 + 5 = 9 > 6$ $\checkmark$
    \item $4 + 6 = 10 > 5$ $\checkmark$
    \item $5 + 6 = 11 > 4$ $\checkmark$
    \item Conclusion: Valid triangle
\end{itemize}

\textbf{Check 3: Scalene Requirement}
\begin{itemize}
    \item $4 \neq 5 \neq 6 \neq 4$ $\checkmark$
    \item Conclusion: Properly scalene
\end{itemize}

\textbf{Check 4: Box Realizability}
\begin{itemize}
    \item $p^2 = \frac{16+25-36}{2} = \frac{5}{2}$ (positive) $\checkmark$
    \item $q^2 = \frac{16-25+36}{2} = \frac{27}{2}$ (positive) $\checkmark$
    \item $r^2 = \frac{-16+25+36}{2} = \frac{45}{2}$ (positive) $\checkmark$
    \item Conclusion: Disphenoid is realizable
\end{itemize}

\textbf{Number Theory Verification:}

\textbf{Check 5: Heron's Formula Calculation}
\begin{align*}
s &= \frac{4+5+6}{2} = \frac{15}{2} \\
A &= \sqrt{s(s-a)(s-b)(s-c)} \\
&= \sqrt{\frac{15}{2} \cdot \frac{7}{2} \cdot \frac{5}{2} \cdot \frac{3}{2}} \\
&= \sqrt{\frac{15 \cdot 7 \cdot 5 \cdot 3}{16}} \\
&= \sqrt{\frac{1575}{16}} \\
&= \sqrt{\frac{225 \cdot 7}{16}} \\
&= \frac{15\sqrt{7}}{4}
\end{align*}
\textbf{Check 6: Total Surface Area}
\begin{align*}
\text{Total} &= 4A = 4 \cdot \frac{15\sqrt{7}}{4} = 15\sqrt{7}
\end{align*}

\textbf{Check 7: Answer Choice Validation}
\begin{itemize}
    \item Our answer: $15\sqrt{7}$
    \item Choice (D): $15\sqrt{7}$ $\checkmark$
    \item Match confirmed
\end{itemize}

\textbf{Optimization Verification:}

\textbf{Check 8: Smaller Triangle Verification}
\begin{itemize}
    \item (2,3,4): $4 + 9 = 13 < 16$ (obtuse) $\times$
    \item (2,4,5): $4 + 16 = 20 < 25$ (obtuse) $\times$
    \item (3,4,5): $9 + 16 = 25$ (right angle) $\times$
    \item (4,5,6): Acute $\checkmark$
    \item Conclusion: (4,5,6) is indeed minimum
\end{itemize}

\subsection*{C. Confidence Calibration}

\textbf{Pre-Verification Confidence: 75\%}
\begin{itemize}
    \item Strong: Systematic approach and clear calculation
    \item Weak: Novel geometric constraint (acute requirement)
    \item Concern: Might have missed smaller valid triangle
\end{itemize}

\textbf{Post-Verification Confidence: 95\%}
\begin{itemize}
    \item Confirmed: (4,5,6) satisfies all constraints
    \item Confirmed: Calculation is correct
    \item Confirmed: No smaller valid triangle exists
    \item Remaining 5\%: Arithmetic errors possible
\end{itemize}

\textbf{Final Decision: SELECT ANSWER (D)}

\end{verificationbox}

\section{Solution Roadmap}

\begin{solutionbox}

\subsection*{Complete Solution Path}

\textbf{Phase 1: Understanding the Problem (1-2 minutes)}

\textbf{Step 1.1: Define Key Terms}
\begin{itemize}
    \item Disphenoid: Tetrahedron with 4 congruent faces
    \item Scalene: All three sides different
    \item Integer side lengths: $a, b, c \in \mathbb{Z}^+$
    \item Goal: Minimize total surface area = $4A$
\end{itemize}

\textbf{Step 1.2: Identify Structure}
\begin{itemize}
    \item Key property: In disphenoid with congruent faces, opposite edges are equal
    \item This means: If faces have sides $a, b, c$, then edges pair up
    \item Implication: Disphenoid can be inscribed in rectangular box
\end{itemize}

\textbf{Checkpoint 1:} Do I understand what makes a valid disphenoid?

\textbf{Phase 2: Discovering the Constraint (2-3 minutes)}

\textbf{Step 2.1: Box Embedding}
\begin{itemize}
    \item Place disphenoid in box with dimensions $p, q, r$
    \item Triangle sides are box face diagonals:
    \begin{align*}
    a^2 &= p^2 + q^2 \\
    b^2 &= p^2 + r^2 \\
    c^2 &= q^2 + r^2
    \end{align*}
\end{itemize}

\textbf{Step 2.2: Solve for Box Dimensions}
\begin{itemize}
    \item Adding equations strategically:
    \begin{align*}
    p^2 &= \frac{a^2 + b^2 - c^2}{2} \\
    q^2 &= \frac{a^2 - b^2 + c^2}{2} \\
    r^2 &= \frac{-a^2 + b^2 + c^2}{2}
    \end{align*}
\end{itemize}

\textbf{Step 2.3: Derive Acute Constraint}
\begin{itemize}
    \item For $p, q, r$ to be real and positive, need:
    \begin{align*}
    a^2 + b^2 - c^2 &> 0 \implies a^2 + b^2 > c^2 \\
    a^2 - b^2 + c^2 &> 0 \implies a^2 + c^2 > b^2 \\
    -a^2 + b^2 + c^2 &> 0 \implies b^2 + c^2 > a^2
    \end{align*}
    \item This is precisely the acute triangle condition!
\end{itemize}

\textbf{Checkpoint 2:} Have I discovered that faces must be acute triangles?

\textbf{Phase 3: Finding Candidate Triangles (1-2 minutes)}

\textbf{Step 3.1: List Small Scalene Triangles}
\begin{enumerate}
    \item (2,3,4): Test $4 + 9 = 13 < 16$ $\to$ obtuse $\times$
    \item (2,4,5): Test $4 + 16 = 20 < 25$ $\to$ obtuse $\times$
    \item (3,4,5): Test $9 + 16 = 25$ $\to$ right angle $\times$
    \item (4,5,6): Test $16 + 25 = 41 > 36$ $\to$ acute $\checkmark$
\end{enumerate}

\textbf{Step 3.2: Verify All Conditions for (4,5,6)}
\begin{itemize}
    \item Scalene: $4 \neq 5 \neq 6$ $\checkmark$
    \item Integer sides: Yes $\checkmark$
    \item Acute: All three inequalities satisfied $\checkmark$
    \item Valid triangle: Triangle inequality holds $\checkmark$
\end{itemize}

\textbf{Checkpoint 3:} Have I found the smallest valid triangle?

\textbf{Phase 4: Calculate Area (2-3 minutes)}

\textbf{Step 4.1: Apply Heron's Formula}
\begin{itemize}
    \item Semiperimeter: $s = \frac{4+5+6}{2} = \frac{15}{2}$
    \item Factors:
    \begin{align*}
    s - a &= \frac{15}{2} - 4 = \frac{7}{2} \\
    s - b &= \frac{15}{2} - 5 = \frac{5}{2} \\
    s - c &= \frac{15}{2} - 6 = \frac{3}{2}
    \end{align*}
\end{itemize}

\textbf{Step 4.2: Compute Area}
\begin{align*}
A &= \sqrt{s(s-a)(s-b)(s-c)} \\
&= \sqrt{\frac{15}{2} \cdot \frac{7}{2} \cdot \frac{5}{2} \cdot \frac{3}{2}} \\
&= \sqrt{\frac{15 \cdot 7 \cdot 5 \cdot 3}{16}}
\end{align*}

\textbf{Step 4.3: Simplify}
\begin{align*}
A &= \sqrt{\frac{15 \cdot 7 \cdot 5 \cdot 3}{16}} \\
&= \sqrt{\frac{15 \cdot 105}{16}} \\
&= \sqrt{\frac{1575}{16}} \\
&= \sqrt{\frac{225 \cdot 7}{16}} \\
&= \frac{15\sqrt{7}}{4}
\end{align*}

\textbf{Checkpoint 4:} Is my calculation correct?

\textbf{Phase 5: Final Answer (30 seconds)}

\textbf{Step 5.1: Total Surface Area}
\begin{align*}
\text{Total Surface Area} &= 4A = 4 \cdot \frac{15\sqrt{7}}{4} = 15\sqrt{7}
\end{align*}

\textbf{Step 5.2: Match to Answer Choices}
\begin{itemize}
    \item Our answer: $15\sqrt{7}$
    \item Answer choice: $\textbf{(D) } 15\sqrt{7}$
\end{itemize}

\subsection*{Alternative Approaches}

\textbf{Alternative 1: Volume Formula Method}
\begin{itemize}
    \item Use disphenoid volume formula:
    \[72V^2 = (-a^2+b^2+c^2)(a^2-b^2+c^2)(a^2+b^2-c^2)\]
    \item For positive volume, all three factors must be positive
    \item This directly gives acute triangle condition
    \item Then proceed with same enumeration
\end{itemize}

\textbf{Alternative 2: Answer Choice Back-Solving}
\begin{itemize}
    \item Calculate area for each answer choice
    \item Divide by 4 to get single face area
    \item Use inverse Heron to find possible side lengths
    \item Check which gives smallest integer sides
    \item More time-consuming but valid
\end{itemize}

\subsection*{Common Pitfalls}

\textbf{Pitfall 1: Testing Non-Acute Triangles}
\begin{itemize}
    \item Error: Calculating area for (2,3,4) or (3,4,5)
    \item Why wrong: These don't form valid disphenoids
    \item Prevention: Always check acute condition first
\end{itemize}

\textbf{Pitfall 2: Arithmetic Errors in Heron's Formula}
\begin{itemize}
    \item Common: Forgetting to take square root, or error in simplification
    \item Prevention: Double-check each step, especially simplifying radicals
    \item Verification: Use calculator if allowed
\end{itemize}

\textbf{Pitfall 3: Forgetting to Multiply by 4}
\begin{itemize}
    \item Error: Giving $\frac{15\sqrt{7}}{4}$ as final answer
    \item Why wrong: This is area of ONE face, not total
    \item Prevention: Re-read problem asking for ``total surface area''
\end{itemize}

\textbf{Pitfall 4: Assuming Smaller Sides = Smaller Area}
\begin{itemize}
    \item Error: Not checking if smaller triangles are valid
    \item Why wrong: Geometric constraints limit options
    \item Prevention: Systematic checking with constraint verification
\end{itemize}

\end{solutionbox}

\section{Competition Strategy Integration}

\begin{competitionbox}

\subsection*{A. Time Management}

\textbf{Time Budget for Problem 24:}
\begin{itemize}
    \item Total allocated: 6-8 minutes (late-band problem)
    \item Phase 1 (Understanding): 1-2 minutes
    \item Phase 2 (Finding constraint): 2-3 minutes
    \item Phase 3 (Testing triangles): 1-2 minutes
    \item Phase 4 (Calculation): 2-3 minutes
    \item Phase 5 (Verification): 1-2 minutes
\end{itemize}

\textbf{Time Checkpoints:}
\begin{itemize}
    \item After 2 minutes: Should understand disphenoid and box structure
    \item After 4 minutes: Should have discovered acute constraint
    \item After 6 minutes: Should have calculated area
    \item After 8 minutes: If not done, consider strategic skip
\end{itemize}

\subsection*{B. Strategic Positioning}

\textbf{When to Attempt:}
\begin{itemize}
    \item Strong geometry background: Attempt after P1-23
    \item Weak geometry background: Skip and use time elsewhere
    \item Medium geometry: Attempt if $>$10 minutes remaining
\end{itemize}

\textbf{Problem Sequencing:}
\begin{itemize}
    \item This is P24 - naturally comes near end
    \item Consider: Is P25 easier for you? Look at it first
    \item Strategy: Do P1-23, assess time, then attempt P24-25
\end{itemize}

\subsection*{C. Risk Assessment}

\textbf{Success Probability:}
\begin{itemize}
    \item With strong geometry: 60-70\% (if acute constraint known)
    \item With standard geometry: 30-40\% (constraint is hard to discover)
    \item With weak geometry: 10-15\% (multiple obstacles)
\end{itemize}

\textbf{Risk Level: HIGH}
\begin{itemize}
    \item Difficulty: Very high (P24/25)
    \item Complexity: Multi-step with subtle constraint
    \item Time cost: 6-8 minutes
    \item Payoff: Same 6 points as easier problems
\end{itemize}

\textbf{Risk Mitigation:}
\begin{itemize}
    \item Quick check: Do I know about disphenoids? (If no: risky)
    \item Quick check: Am I comfortable with Heron's formula? (If no: risky)
    \item Quick check: Do I have $>$8 minutes? (If no: skip)
    \item Backup: If stuck after 3 minutes, can I guess strategically?
\end{itemize}

\subsection*{D. Abandonment Criteria}

\textbf{When to Abandon:}

\textbf{Scenario 1: After 3 minutes, still don't understand disphenoid}
\begin{itemize}
    \item Action: Skip to P25 or return to easier unsolved problems
    \item Reason: Fundamental misunderstanding, unlikely to solve
\end{itemize}

\textbf{Scenario 2: After 5 minutes, calculation errors persist}
\begin{itemize}
    \item Action: Make best guess from answer choices
    \item Reason: Diminishing returns, time better spent elsewhere
\end{itemize}

\textbf{Scenario 3: Less than 5 minutes remaining in exam}
\begin{itemize}
    \item Action: Quick guess or strategic elimination
    \item Reason: Not enough time for careful calculation
\end{itemize}

\textbf{Strategic Guessing (if must abandon):}
\begin{itemize}
    \item Eliminate (A) $\sqrt{3}$: Too small
    \item Eliminate (C) $15$: Too clean for scalene requirement
    \item Guess between (B), (D), (E)
    \item Slight preference for (D): Form matches expected structure
\end{itemize}

\subsection*{E. Answer Recording}

\textbf{Recording Protocol:}
\begin{itemize}
    \item Write answer clearly: D
    \item Double-check bubble: (D) is filled
    \item Cross-reference: Problem 24 = bubble 24
    \item Final check: No erasure marks causing ambiguity
\end{itemize}

\subsection*{F. Post-Problem Strategy}

\textbf{If Solved Successfully:}
\begin{itemize}
    \item Confidence boost for P25
    \item Time check: How much left?
    \item Review: Any earlier problems to revisit?
\end{itemize}

\textbf{If Abandoned:}
\begin{itemize}
    \item No regret: P24 is very hard
    \item Focus: Ensure all easier problems are done
    \item Time use: Better spent on achievable points
\end{itemize}

\end{competitionbox}

\section{Summary and Key Takeaways}

\begin{keyinsight}
\textbf{Critical Insights for This Problem:}

\begin{enumerate}
    \item The defining property of a disphenoid (all faces congruent) forces opposite edges to be equal, which leads to the box-embedding structure.
    
    \item The geometric realizability constraint (box dimensions must be real and positive) translates to an algebraic condition that the face triangle must be \textit{acute}.
    
    \item The acute triangle requirement eliminates many small integer triangles: (2,3,4), (2,4,5), and (3,4,5) all fail.
    
    \item The smallest scalene acute triangle with integer sides is (4,5,6), which gives area $\frac{15\sqrt{7}}{4}$ per face.
    
    \item Total surface area = $4 \times \frac{15\sqrt{7}}{4} = 15\sqrt{7}$.
\end{enumerate}
\end{keyinsight}

\begin{protip}
\textbf{Generalization:} For any triangle $(x-1, x, x+1)$ (consecutive integers), the area by Heron's formula is:
\[A = \frac{x\sqrt{3(x^2-4)}}{4}\]
For $x=5$ (giving 4-5-6), this yields $A = \frac{5\sqrt{3 \cdot 21}}{4} = \frac{5\sqrt{63}}{4} = \frac{15\sqrt{7}}{4}$.

This pattern is useful for quickly recognizing consecutive integer triangles!
\end{protip}

\section{Final Answer}

\begin{center}
\Large
\fbox{\textbf{Answer: (D) } $15\sqrt{7}$}
\end{center}

\textbf{Solution Summary:}
\begin{itemize}
    \item A disphenoid with scalene faces requires those faces to be acute triangles
    \item The smallest scalene acute triangle with integer sides is (4, 5, 6)
    \item Using Heron's formula with semiperimeter $s = \frac{15}{2}$:
    \[A = \sqrt{\frac{15}{2} \cdot \frac{7}{2} \cdot \frac{5}{2} \cdot \frac{3}{2}} = \frac{15\sqrt{7}}{4}\]
    \item Total surface area = $4A = 15\sqrt{7}$
\end{itemize}

\end{document}